% Generated by analyze_schedule_ratio_sensitivity.py; do not hand-edit.
\begin{table}[H]
\centering\small
\begin{tabular}{@{}lcc@{}}
\toprule
Policy & \shortstack{True-boundary-exceeding\\final recommendation (\%)} & \shortstack{Above-boundary simulated\\assignments (count/40)} \\
\midrule
pure cEI & $11.0{\pm}2.4$ & $6.82{\pm}0.47$ \\
cKG & $18.5{\pm}2.8$ & $7.40{\pm}0.35$ \\
pure tMSE & $13.0{\pm}2.6$ & $11.54{\pm}0.51$ \\
cEI--tMSE & $6.0{\pm}1.6$ & $9.71{\pm}0.41$ \\
cEI--cEI--tMSE cycle & $15.0{\pm}2.6$ & $8.98{\pm}0.42$ \\
\bottomrule
\end{tabular}
\caption{Exploratory post hoc schedule-composition sensitivity for the full-panel OSA latent-gate cell at $\tau=0.7$. Entries are means $\pm$ Monte Carlo SE computed across 100 independent matched replicate sets after averaging the two strata within each set (200 simulated trials per policy). The cEI--tMSE is the study-defined reference schedule. Under this full-panel design, $N_{\max}=40$, four initialization assignments, and cohorts of two give 18 scheduled adaptive-cohort positions: nine cEI and nine tMSE. An empty eligible gate can instead invoke the common greatest-feasibility fallback. The cEI--cEI--tMSE cycle gives 12 cEI and six tMSE adaptive cohorts. It was specified after the original study results were examined and was not tuned. Assignment counts include all 40 simulated assignments. This analysis is descriptive and nonconfirmatory.}
\label{tab:schedule_ratio_posthoc}
\end{table}
